\documentclass[11pt]{article}
\usepackage{amsmath,hyperref}
\title{A delivery port can reopen before the old return settles}
\author{}
\date{}
\begin{document}
\maketitle

In a two-carrier warehouse simulation, clearing the delivery port lets the next order reach the port at tick 5 rather than 16. The earlier carrier still owes its return at tick 12. But the same early view of an empty port can precede an old carrier's authorized re-entry; in that case the next delivery cannot finish before tick 10. An ordinary controller that waits only for its current movement grant and the port to be clear reaches the same optimum. The result identifies what a controller must know to \emph{promise} early service, not a faster planner.

\paragraph{Native effects and the operating rule.}
This replay executes the exact selected agent, shelf, movement, grid and step methods of \href{https://github.com/semitable/robotic-warehouse/blob/96fbc64e3eae5fee915e0d390f864fa06ddccd47/rware/warehouse.py}{Robotic Warehouse at 96fbc64}. It uses the one-mover headless harness from the \href{https://github.com/eternal-november/eos-research-programmes/pull/418}{previous source-and-return replay}. There is no Gym environment or network planner in this replay. Both carriers run under an authored shared budget of at most one non-idle operation per tick; this is not a limitation of the full simulator. In native two-stage reward, the first delivery changes the sole queued request and gives half credit. The other half is paid only when the delivering carrier drops its shelf off the highway. The mission being scored consists of two specified deliveries and their returns, not later requests created by the simulator.

The old carrier initially delivers from below and then leaves north at tick 1. In the no-re-entry scenario it eventually parks its original shelf north of the port. In the other scenario, it holds irrevocable scheduled movement to turn, re-enter, and leave to the south. In the paired control it re-enters at tick 5, leaves at 9, and drops at 12. Both controls schedule the actual drop at 12. The second shelf awaits at cell $(1,1)$, with the next carrier facing east toward the port at $(4,1)$. A reservation for an old action is not one that the next carrier can reschedule or spend. All old actions execute through native mechanics. A reservation for the next carrier also cannot pre-empt an old occupied cell. The operator requires continued service of the old commitment, legal custody of both shelves, and the earliest next delivery possible under those constraints.

\paragraph{The exact cutoff.}
Count ticks not reserved for the old carrier. The next carrier can pick up and move east twice at the first three such grants. These intermediate cells are not on the old route. The earliest delivery is then the first unreserved tick with the port empty at its start. Necessity follows because pickup and at least three native moves are required, the old reservation excludes using that tick, and entering an occupied port cannot deliver. Sufficiency follows by staging on the last cell before the port, checking the current grant, and advancing when empty. Because this route never intercepts an old action, independent pre-delivery work cannot postpone the old return. The controller needs no forecast to execute this earliest policy. A controller promising a deadline at the cut after tick 1 must nevertheless allow every old schedule compatible with that cut. A forecast that the port is currently empty cannot certify that it will stay empty or that the remaining moves have unreserved slots. With precisely the admissible pair in this replay and no further information at the promise cut, the sharp robust commitment is tick 10; changing future entitlement or narrowing the possible schedules can change it. Merely waiting until the old carrier proves its \emph{entire} return has settled misses the earlier legal opportunity.

\paragraph{Witnesses and comparisons.}
The two scenarios have the same native observation after tick 1, including shelf identities, delivery flag, and queue. The direct native search, the factored computation, the exact cutoff, and one feedback policy using only current grant and port occupancy agree on the earliest delivery in all 50 developed scenarios, including the globally paused comparison. The direct reference explores 13,099 native transitions with actual old reservations. Replaying an uncorrected fixed route based on the early empty-port view makes the native gate reject the move into the port at tick 8 in the re-entry case; no collision is claimed. Separate continuation traces finish both returns as well. In the no-re-entry comparison the scoped policy finishes both obligations at tick 12 against 19 under pause; in the re-entry comparison it finishes at 14 against 19. The old drop remains at 12. The latter overall finish times are attained schedules, not claimed to be optimal for unrestricted return planning.

For \href{https://github.com/eternal-november/eos-research-programmes/pull/419}{the fence-and-drain problem}, this shows why a useful release certificate can be specific to a port and the outstanding schedule rather than a certificate of global settlement. It does not produce a distributed certificate of future commitments. The plan comparison assumes reservations are exogenous, obeyed, and surfaced as current grants; changing the source or allowing reactive old moves requires checking those assumptions again. The old carrier's contingent re-entry is an authored native-source development case, not an independently selected fault, source validation of the grant system, or a positive advantage over the strongest matched policy. No I1 or R2 acceptance is claimed.
\end{document}
